\documentclass[10pt]{article}

% ---------- page layout ----------
\usepackage[landscape, margin=0.3cm]{geometry}

% ---------- fonts / encoding ----------
\usepackage[T1]{fontenc}
\usepackage{lmodern}

% ---------- tables ----------
\usepackage{longtable}
\usepackage{booktabs}
\usepackage{array}
\usepackage{makecell}          % \makecell for line breaks inside cells
\usepackage{siunitx}           % consistent unit formatting

% ---------- misc ----------
\usepackage[colorlinks=true,linkcolor=blue,urlcolor=blue,citecolor=blue]{hyperref}
\usepackage[table]{xcolor}

% ---------- compact spacing ----------
\renewcommand{\arraystretch}{1.25}
\setlength{\tabcolsep}{4pt}

% ---------- convenience macros ----------
\newcommand{\um}{\textmu{}m}
\newcommand{\srcref}[4]{\href{https://doi.org/#4}{\textcolor{blue}{#1,\penalty0\,#2,\penalty0\,#3}}}
\newcommand{\srcurl}[3]{\href{#3}{\textcolor{blue}{#1,\penalty0\,#2}}}

\pagestyle{empty}

\begin{document}

% ======================================================================
%  Beam Source Comparison
% ======================================================================
\begin{center}
{\Large\bfseries Beam Source Parameters for Imaging}\\[4pt]
{\small Photon flux values are typical for an undulator imaging beamline (0.1\% BW, at sample).
Synchrotron values are beamline-dependent; representative numbers are given.
Ph/pulse is derived from total flux $\div$ bunch rate.}
\end{center}

\vspace{4pt}

\scriptsize
\renewcommand{\arraystretch}{1.20}
\setlength{\tabcolsep}{3pt}

\begin{longtable}{%
  l                                               % 1  Facility
  l                                               % 2  Type
  >{$}l<{$}                                       % 3  Particle & Energy
  >{$}l<{$}                                       % 4  Total flux at sample
  >{$}l<{$}                                       % 5  Rep rate / bunch spacing
  >{$}l<{$}                                       % 6  Pulse duration
  >{$}l<{$}                                       % 7  Particles per pulse
  >\raggedright\arraybackslash p{3.0cm}            % 8  Notes
  >\raggedright\arraybackslash p{2.8cm}            % 9  Source
}

\toprule
\textbf{Facility}
  & \textbf{Type}
  & \textbf{Particle \& Energy}
  & \makecell[tl]{\textbf{Total flux}\\\textbf{(ph/s or e$^-$/s)}}
  & \makecell[tl]{\textbf{Rep.\ rate}\\\textbf{(bunch spacing)}}
  & \makecell[tl]{\textbf{Pulse}\\\textbf{duration}}
  & \makecell[tl]{\textbf{Particles}\\\textbf{per pulse}}
  & \textbf{Notes}
  & \textbf{Source} \\
\midrule
\endfirsthead

\toprule
\textbf{Facility}
  & \textbf{Type}
  & \textbf{Particle \& Energy}
  & \makecell[tl]{\textbf{Total flux}\\\textbf{(ph/s or e$^-$/s)}}
  & \makecell[tl]{\textbf{Rep.\ rate}\\\textbf{(bunch spacing)}}
  & \makecell[tl]{\textbf{Pulse}\\\textbf{duration}}
  & \makecell[tl]{\textbf{Particles}\\\textbf{per pulse}}
  & \textbf{Notes}
  & \textbf{Source} \\
\midrule
\endhead

\bottomrule
\endlastfoot

%% ===== PETRA III =====
\makecell[tl]{PETRA III\\(DESY)}
  & \makecell[tl]{3\textsuperscript{rd} gen.\\synchrotron}
  & \gamma,\; 5\text{--}200\,\text{keV}
  & {\sim}\,10^{12}\text{--}10^{13}
  & \makecell[tl]{$62.5\,\text{MHz}$ (16\,ns)\\or $5.2\,\text{MHz}$ (192\,ns)}
  & {\sim}\,100\,\text{ps}
  & {\sim}\,10^{4}\text{--}10^{5}
  & \makecell[tl]{120\,mA continuous\\or 100\,mA timing mode.\\6\,GeV, 2304\,m ring.\\960 or 40 bunches}
  & \makecell[tl]{\srcurl{DESY}{Facility info}{https://photon-science.desy.de/facilities/petra_iii/facility_information/}} \\
\addlinespace

%% ===== PETRA IV =====
\makecell[tl]{PETRA IV\\(DESY)}
  & \makecell[tl]{4\textsuperscript{th} gen.\\synchrotron}
  & \gamma,\; 5\text{--}200\,\text{keV}
  & {\sim}\,10^{12}\text{--}10^{13}
  & \text{same as PETRA III}
  & {\sim}\,30\text{--}100\,\text{ps}
  & {\sim}\,10^{4}\text{--}10^{5}
  & \makecell[tl]{100--1000$\times$ brilliance\\vs PETRA III.\\Total flux similar, but\\focused into smaller spot\\$\Rightarrow$ higher fluence.\\$\varepsilon \approx 20$\,pm\,rad}
  & \makecell[tl]{\srcref{2018}{Schroer}{J.\,Synch.\,Rad.}{10.1107/S1600577518007907};\\
  \srcref{2022}{Schroer}{EPJ Plus}{10.1140/epjp/s13360-022-03517-6}} \\
\addlinespace

%% ===== ESRF-EBS =====
\makecell[tl]{ESRF-EBS\\(Grenoble)}
  & \makecell[tl]{4\textsuperscript{th} gen.\\synchrotron}
  & \gamma,\; 6\text{--}250\,\text{keV}
  & {\sim}\,10^{13}
  & \makecell[tl]{$352\,\text{MHz}$ (2.8\,ns)\\or $5.7\,\text{MHz}$ (176\,ns)}
  & {\sim}\,50\text{--}200\,\text{ps}
  & {\sim}\,10^{4}\text{--}10^{5}
  & \makecell[tl]{200\,mA, 6\,GeV,\\844\,m ring.\\Uniform (992 bunches)\\or 16-bunch mode.\\$\varepsilon \approx 130$\,pm\,rad}
  & \makecell[tl]{\srcref{2023}{Raimondi}{Comm.\,Phys.}{10.1038/s42005-023-01195-z};\\
  \srcurl{ESRF}{Parameters}{https://www.esrf.fr/home/UsersAndScience/Accelerators/parameters.html}} \\
\addlinespace

%% ===== EuXFEL SP =====
\makecell[tl]{EuXFEL\\SP mode\\(DESY)}
  & \makecell[tl]{XFEL\\(burst)}
  & \gamma,\; 3\text{--}25\,\text{keV}
  & {\sim}\,2.7 \times 10^{16}\,\text{(peak)}
  & \makecell[tl]{$4.5\,\text{MHz}$ (222\,ns)\\intra-train;\\10\,Hz trains\\(up to 2700/train)}
  & {\sim}\,10\text{--}100\,\text{fs}
  & {\sim}\,10^{12}
  & \makecell[tl]{17.5\,GeV.\\0.65\,ms trains at 10\,Hz.\\Up to 27\,k pulses/s.\\${\sim}1.5$\,mJ/pulse.\\``Diffract before destroy''}
  & \makecell[tl]{\srcref{2020}{Sobolev}{Comm.\,Phys.}{10.1038/s42005-020-0362-y};\\
  \srcref{2014}{Brinkmann}{NIMA}{10.1016/j.nima.2014.09.033}} \\
\addlinespace

%% ===== EuXFEL CW =====
\makecell[tl]{EuXFEL\\CW mode\\(DESY)}
  & \makecell[tl]{XFEL\\(CW, planned)}
  & \gamma,\; \text{soft--hard}
  & {\sim}\,1.2 \times 10^{11}\,\text{ph/mm}^2\text{/s}
  & 150\,\text{kHz}\;(6.7\,\text{\textmu s})
  & {\sim}\,\text{fs}
  & {\sim}\,10^{4}\,\text{ph/px}\;\text{(est.)}
  & \makecell[tl]{Max 7.8\,GeV.\\0.5\,nC/bunch,\\0.125\,mA average.\\Continuous, no trains.\\CoRDIA design target:\\$1.5\times10^9$\,ph/px/s\\@110\,\um{} pitch}
  & \makecell[tl]{\srcref{2014}{Brinkmann}{NIMA}{10.1016/j.nima.2014.09.033};\\
  \srcref{2025}{Trunk}{JPCS}{10.1088/1742-6596/3010/1/012141}} \\
\addlinespace

%% ===== Talos F200X =====
\makecell[tl]{Talos F200X\\(Thermo Fisher)}
  & S/TEM
  & e^{-},\; 200\,\text{keV}
  & \makecell[tl]{$9.4 \times 10^{9}$\\(1.5\,nA)}
  & \text{DC (continuous)}
  & \text{DC}
  & \text{N/A (DC beam)}
  & \makecell[tl]{X-FEG.\\Probe: 1.5\,nA @1\,nm}
  & \makecell[tl]{\srcurl{TF}{Talos F200X DS}{https://assets.thermofisher.com/TFS-Assets/MSD/Datasheets/Talos-F200X-Datasheet.pdf}} \\
\addlinespace

%% ===== Talos F200 =====
\makecell[tl]{Talos F200\\(Thermo Fisher)}
  & S/TEM
  & e^{-},\; 200\,\text{keV}
  & \makecell[tl]{$1.4 \times 10^{9}$\\(220\,pA)}
  & \text{DC (continuous)}
  & \text{DC}
  & \text{N/A (DC beam)}
  & \makecell[tl]{X-CFEG.\\Probe: 220\,pA (EELS)}
  & \makecell[tl]{\srcurl{TF}{Talos F200 DS}{https://www.thermofisher.com/us/en/home/electron-microscopy/products/transmission-electron-microscopes.html}} \\
\addlinespace

%% ===== Spectra 200/300 =====
\makecell[tl]{Spectra 200/300\\(Thermo Fisher)}
  & S/TEM
  & e^{-},\; 200\text{--}300\,\text{keV}
  & \makecell[tl]{$6.2 \times 10^{9}$\\(${>}1$\,nA)}
  & \text{DC (continuous)}
  & \text{DC}
  & \text{N/A (DC beam)}
  & \makecell[tl]{X-CFEG + S-CORR.\\${>}1$\,nA @${<}76$\,pm;\\100\,pA @50\,pm;\\30\,pA @50\,pm (mono)}
  & \makecell[tl]{\srcurl{TF}{Spectra DS}{https://www.thermofisher.com/us/en/home/electron-microscopy/products/transmission-electron-microscopes.html}} \\
\addlinespace

%% ===== Titan Themis =====
\makecell[tl]{Titan Themis\\(Thermo Fisher)}
  & S/TEM
  & e^{-},\; 300\,\text{keV}
  & \makecell[tl]{$1.2 \times 10^{10}$\\(2\,nA)}
  & \text{DC (continuous)}
  & \text{DC}
  & \text{N/A (DC beam)}
  & \makecell[tl]{X-FEG, Cs-corr.\\Probe: 2\,nA @0.2\,nm}
  & \makecell[tl]{\srcurl{TF}{Titan Themis DS}{https://www.thermofisher.com/us/en/home/electron-microscopy/products/transmission-electron-microscopes.html}} \\
\addlinespace

%% ===== JEOL ARM300F2 =====
\makecell[tl]{JEM-ARM300F2\\(JEOL)}
  & S/TEM
  & e^{-},\; 300\,\text{keV}
  & \makecell[tl]{$3.1 \times 10^{6}$\\(0.5\,pA)}
  & \text{DC (continuous)}
  & \text{DC}
  & \text{N/A (DC beam)}
  & \makecell[tl]{Cold FEG.\\Probe: 0.5\,pA\\(0.15\,pA low-dose)}
  & \makecell[tl]{\srcurl{JEOL}{ARM300F2}{https://www.jeol.com/products/scientific/tem/JEM-ARM300F2.php}} \\
\addlinespace

%% ===== JEOL ARM200F =====
\makecell[tl]{JEM-ARM200F\\(JEOL)}
  & S/TEM
  & e^{-},\; 200\,\text{keV}
  & \makecell[tl]{\text{probe dep.}}
  & \text{DC (continuous)}
  & \text{DC}
  & \text{N/A (DC beam)}
  & \makecell[tl]{Cs-corrected.\\Current density:\\8\,nA/nm$^2$ in\\focused probe}
  & \makecell[tl]{\srcurl{JEOL}{ARM200F}{https://www.jeol.com/products/scientific/tem/JEM-ARM200F_NEOARM.php}} \\
\addlinespace

%% ===== ELSA =====
\makecell[tl]{ELSA X3ED\\(Uni Bonn)}
  & \makecell[tl]{Test beam\\(extracted $e^-$)}
  & e^{-},\; 0.8\text{--}3.2\,\text{GeV}
  & \makecell[tl]{$300\,\text{Hz}\text{--}4\times10^{8}$\\(aA--100\,pA)}
  & \text{DC (continuous)}
  & \text{DC}
  & \text{N/A (DC beam)}
  & \makecell[tl]{Slow resonance\\extraction from\\stretcher ring.\\1--10\,mm beam spot}
  & \makecell[tl]{\srcurl{ELSA}{Facility}{https://www.pi.uni-bonn.de/elsa/en}} \\

\end{longtable}

\vspace{4pt}
\noindent\textbf{Notes:}
\begin{itemize}\setlength{\itemsep}{1pt}
  \item \textbf{Ph/pulse} for synchrotrons is derived from total beamline flux $\div$ bunch repetition rate.
  \item \textbf{PETRA IV} brilliance increase is 100--1000$\times$ vs PETRA III, but \emph{total} flux
    (integrated over the full beam) is similar; the gain comes from concentrating photons into a smaller phase space.
  \item \textbf{EuXFEL SP} peak flux is computed as $10^{12}$\,ph/pulse $\times$ 27\,k\,pulses/s;
    time-averaged flux is much lower (0.65\% duty cycle).
  \item \textbf{All electron sources} (S/TEMs, ELSA) deliver DC continuous beams.
    Pulsed operation requires third-party modifications, typically either
    \emph{laser-driven photoemission} (replacing or augmenting the electron gun)
    or \emph{RF cavity beam choppers} (deflecting the beam into timed bunches).
  \item \textbf{S/TEM} total flux is the probe current converted to e$^-$/s.
    Fluence at the detector depends on magnification and illuminated area.
  \item \textbf{ELSA} beam is DC-extracted via slow resonance extraction,
    so ``pulse duration'' and ``particles per pulse'' do not apply.
\end{itemize}

\end{document}
